\abstract % Структурный элемент: РЕФЕРАТ

\keywords{ОНЧЕЙН-ДАННЫЕ, ИНДЕКСАЦИЯ ДАННЫХ, БЛОКЧЕЙН BITCOIN, МОДУЛЬНЫЙ МОНОЛИТ, STATELESS-АРХИТЕКТУРА, RUST, POSTGRESQL, REST API, MEMPOOL, REORG, UTXO}

Выпускная квалификационная работа посвящена разработке системы индексации
ончейн-данных сети Bitcoin. Актуальность темы определяется необходимостью
создания специализированного программного слоя, обеспечивающего
структурированное хранение и прикладную выдачу блокчейн-данных с
предсказуемой производительностью чтения, устойчивостью к изменению
канонической цепи и возможностью интеграции с внешними информационными
системами.

Объектом исследования является процесс получения, нормализации, хранения и
предоставления ончейн-данных сети Bitcoin. Предметом исследования выступают
архитектурные и алгоритмические решения, обеспечивающие индексацию блоков,
транзакций, адресных балансов, UTXO и неподтвержденных транзакций, а также
согласованное предоставление этих данных прикладным потребителям.

Цель работы состоит в проектировании и реализации системы индексации
ончейн-данных, которая получает данные из узла Bitcoin Core по JSON-RPC,
нормализует их, сохраняет в PostgreSQL и предоставляет через REST API,
административный интерфейс и командные средства сопровождения. В ходе работы
были проанализированы особенности предметной области, сформированы требования
к системе, спроектированы архитектура и модель данных, реализованы механизмы
индексации подтвержденной цепи, обработки reorg, синхронизации mempool,
мониторинга и контейнерного развертывания.

Практическим результатом работы является программный комплекс
\texttt{Bitcoin Blockchain Indexer}, включающий серверное приложение на Rust,
схему и миграции базы данных, административную панель, командный интерфейс,
эксплуатационную документацию и интеграционные тесты. Разработанное решение
может использоваться как инфраструктурный компонент аналитических,
исследовательских и сервисных систем, работающих с ончейн-данными Bitcoin.
